\documentclass[12pt, a4paper]{exam}
\usepackage{graphicx}
%\usepackage[left=0.8in, top=0.7in, total={6.2in,8in}]{geometry}
\usepackage[normalem]{ulem}
\renewcommand\ULthickness{1.0pt}   %%---> For changing thickness of underline
\setlength\ULdepth{1.3ex}%\maxdimen ---> For changing depth of underline
\usepackage{amssymb} 

\begin{document}
	%\thispagestyle{empty}
	\noindent
	\begin{minipage}[l]{0.1\textwidth}
		\noindent
		%\includegraphics[width=2.4\textwidth]{uct.png}
	\end{minipage}
\hfill
\begin{minipage}[c]{1.0\textwidth} % use 0.8 when using the image
	\begin{center}
		
		{
		\large \ \ \par
		\large \ \ \par
		\large  Alexander Cristaudo \par
		\small	CRSALE010	\par
		\large	MAM1019H	\par
	\large \textbf{Hand-in 4}	\par
\small	Date: August 26, 2022}
	\end{center}
\end{minipage}
\par
\vspace{0.2in}
\noindent
\uline{ \ \ \ \ \ \ \ \ \ \ \ \ \ \ \ \ \ \ \ \ \ \ \ \ \ \ \ \ \ \ \ \ \ \ \ \ \ \ \ \ \ \ \ \ \ \ \ \ \ \ \ \ \ \ \ \ \ \ \ \ \ \ \ \ \ \ \ \ \ \ \ \ \ \ \ \ \ \  \ \ \ \ \ \ \ \ \ \ \ \ \ \  \ \ \ \ \ \ \ \ \ \ \ \ \ \ \ \ \ \ \ \ \ \ }
\par 
\vspace{0.15in}
\noindent
\centering
{\small \bfseries 	Answers}

\begin{questions}
	\pointsdroppedatright
	\question
	\begin{parts}
		\part[9] False  
		\part[9] True
		\part False
		\part False
	\end{parts}
\vspace{0.2in}
	\pointsdroppedatright
	\question Some $a \in \mathbb{N}-\{1\}$ is minimal if there is no $x \in \mathbb{N}-\{1\}$ such that $x \ne a$ and $x | a$. Now, this means that $p \in \mathbb{N}$, some prime number, is minimal, because there is no $x \in \mathbb{N}-\{1\}$ such that $x | p$ (otherwise it is not a prime number). Now suppose there is some composite number $c$ which is minimal. Then $\exists p \in \mathbb{N}, p $ is a prime number, and $p \ne a$ and $p | a$ so $a$ is not minimal, so the set of all minimal elements is the set of all prime numbers. \\ Now some $a \in \mathbb{N}-\{1\}$ is maximal if there is no $x \in \mathbb{N}-\{1\}$ with $x \ne a$ and $a | x$. Now this means that $0$ is maximal, because there is no $x \in \mathbb{N}-\{1\}$ where $0 | x$. However, this is the only maximal element. To show this, suppose $a \in \mathbb{N}-\{0,1\}$. Now $a | 0$ so $a$ is not maximal, thus $0$ is the only maximal element. 
		\vspace{0.05in}
		\question  To prove this, we have to prove both implications (directions) \\ $(\Longrightarrow)$ Suppose $a \in A$ is the minimum of $A$ under $\le$. This means that $a \le x$ for all $x \in A$. Now $a \le x \iff x \le ^{-1} a$, so $x \le ^{-1} a, \  \forall x \in A$. This is precisely the definition for the maximum of $A$ under $\le ^{-1}$. This means that: \\ $(a$ is the minimum of $A$ under $\le) \Longrightarrow (a$ is the maximum of $A$ under $\le ^{-1})$ \\ 
		
		$(\Longleftarrow)$  Suppose $a \in A$ is the maximum of $A$ under $\le ^{-1}$. This means that $x \le ^{-1} a$ for all $x \in A$. Now $x \le ^{-1} a \iff a \le x$, so $a \le x, $ for all $ x \in A$. This is the definition for the minimum of $A$ under $\le$. Thus, we have that: \\ $(a$ is the maximum of $A$ under $\le ^{-1}) \Longrightarrow (a$ is the minimum of $A$ under $\le)$ \\ 
		Finally, since both directions have been proven, we can conclude that \\ $(a$ is the minimum of $A$ under $\le) \iff (a$ is the maximum of $A$ under $\le ^{-1})$
\end{questions}
\vspace{0.75in}
\end{document}